% Retrieval Augmented Generation (RAG)
A RAG task
is a set of 3-tuples $\Pi=(Q,A,D)$ where
% $S$ is a set of documents,
$Q$ is the question, $A$ is the ground-truth answer, and $D=(D_1,\cdots D_K)$ is
a sequence of $K$ documents
some of which contain information relevant to $Q$ and $A$.
We assume $D$ is given and focus on generating $A$,
while retrieving $D$ from a larger corpus is outside the scope of this paper.
Using a language model $p_\theta(A|Q,D)$,
the goal of the task is to sample an answer $\tilde{A}$ that is as close to $A$ as possible.
% The performance is measured by two metrics:
% the \emph{F1 metric}, which is the token-level Jaccard distance between $A$ and $\tilde{A}$,
% and \emph{TTFT} (Time to First Token), which is the sum of the runtime for
% any additional processing for KV caches and
% the runtime for prefilling the query $Q$.
% TTFT does not include the runtime for prefilling the system prompt.
% Due to the different runtime-accuracy trade-offs,
% proper evaluation requires comparing its Pareto front.

